\section{MVP operations}\label{app:ops}

\Cref{tab:ops} summarizes the MVP operations of the specification's opcode reference: typing,
the domain of each partial operation (its obligation) and the commutative ones. Integer
arithmetic requires one of the modes \kw{checked} (partial: no overflow), \kw{wrap} (modular)
or \kw{sat} (saturating). Every float operation is the IEEE~754-2019 operation rounded to
nearest-even, on the one-NaN domain.

\begin{table}[h]
\centering
\footnotesize
\renewcommand{\arraystretch}{1.15}
\begin{tabularx}{\linewidth}{@{}l>{\raggedright\arraybackslash}X>{\raggedright\arraybackslash}X@{}}
\toprule
Family & Operations & Domain (obligation); commutative operands \\
\midrule
Structure & \op{TUPLE}, \op{PROJ}, \op{SELECT} (both branches evaluated), \op{CAST}, \op{AS\_BYTES}
  (a view), \op{OK}, \op{ERR}, \op{IS\_OK}, \op{GET\_OK}, \op{GET\_ERR}; \op{IDENTITY} is text only
  & \op{CAST}: the value fits when lossy, and a float source is finite; \op{GET\_OK}: an Ok;
  \op{GET\_ERR}: an Err \\
Integer & \op{ADD}, \op{SUB}, \op{MUL}, \op{DIV}, \op{MOD}, \op{NEG}, each with a mode &
  checked: no overflow; \op{DIV}, \op{MOD}: non-zero divisor, and for signed \op{DIV} not
  $\mathrm{MIN}/{-1}$; \op{ADD}, \op{MUL} commutative \\
Bits & \op{AND}, \op{OR}, \op{XOR}, \op{NOT}, \op{SHL}, \op{SHR} (checked or wrap), \op{ROTL},
  \op{ROTR} & checked shifts: amount below the width; \op{AND}, \op{OR}, \op{XOR} commutative \\
Float & \op{FADD}, \op{FSUB}, \op{FMUL}, \op{FDIV}, \op{FMA}, \op{FNEG}, \op{FABS}, \op{FMIN},
  \op{FMAX}, \op{FCLAMP}, \op{FCOPYSIGN}, \op{FSQRT} & \op{FCLAMP}: bounds not NaN and
  $lo\le hi$; \op{FCOPYSIGN}: sign operand not NaN; \op{FADD}, \op{FMUL}, \op{FMIN}, \op{FMAX}
  commutative, \op{FMA} in its first two operands \\
Compare, logic & \op{EQ}, \op{NE}, \op{LT}, \op{LE}, \op{GT}, \op{GE} (mode \kw{s}, \kw{u} or
  \kw{f}), \op{BAND}, \op{BOR}, \op{BNOT} & total; \op{EQ}, \op{NE}, \op{BAND}, \op{BOR}
  commutative \\
Collections & \op{INDEX}, \op{UPDATE}, \op{LEN}, \op{SLICE}, \op{CONCAT} (indices are
  \kw{u64}) & index below the length; $lo\le hi\le$ length; concatenated length below $2^{64}$ \\
Combinators & \op{SIMD\_MAP}, \op{FOLD\_BOUNDED} (with a named invariant), \op{GENERATE},
  \op{ASSERT\_GUARD} (lazy) & see \cref{tab:rules} \\
Composition & \op{CALL\_HASH} & the callee's precondition \\
Effects & \op{FILE\_READ}, \op{FILE\_WRITE}, \op{CLOCK\_READ} & take and return the effect
  token; results come from the environment \\
\bottomrule
\end{tabularx}
\caption{The MVP operations. Partial operations trap in the interpreter outside their domain,
and a verified node never applies one there.}
\label{tab:ops}
\end{table}

\section{The VG-Sym grammar}\label{app:grammar}

The grammar below is reproduced from the specification. A port (\kw{\%12}) and a label
(\kw{@soft\_knee}) are single tokens; every word the grammar spells out is reserved.

\begin{lstlisting}[language={},basicstyle=\ttfamily\scriptsize,frame=lines]
Module     ::= "module" Ident ";" Item*
Item       ::= MeshQuery | NodeDef | SpecDef
MeshQuery  ::= "@MESH_QUERY" Label "{" "signature" ":" FnType ";" "effects" ":" Effects ";"
               ("pre" ":" Formula ";")? "post" ":" Formula ";" ("intent" ":" String ";")? "}"
SpecDef    ::= "SPEC" Label NodeBody          (* executable reference spec *)
NodeDef    ::= "NODE" Label Attrs? NodeBody
NodeBody   ::= "{" "SIGNATURE" ":" Params "->" Results ";"
                   "EFFECTS" ":" Effects ";"
                   "CONTRACT" "{" Clause* "}"
                   Learnable?
                   "BODY" "{" Stmt+ "}" "}"
Attrs      ::= "[" Attr ("," Attr)* "]"
Attr       ::= "intent" ":" String | "export" | "inline"
Params     ::= "(" (Ident ":" Type ("," Ident ":" Type)*)? ")"
Results    ::= "(" Type ("," Type)* ")"
FnType     ::= "(" (Type ("," Type)*)? ")" "->" Results
Clause     ::= "PRE" ":" Formula ";" | "POST" ":" Formula ";"
             | "INVARIANT" Label "(" Ident "," Ident ")" ":" Formula ";"
             | "DECREASES" ":" Formula ";" | "COST" ":" Formula ";" | "BUDGET" ":" Formula ";"
Learnable  ::= "LEARNABLE" "{" (Ident ":" FloatType "in" "[" Signed "," Signed "]"
               ("loss" ":" Label)? ";")+ "}"
Stmt       ::= Binding ("," Binding)* "=" Op "(" Args? ")" ";"
             | "RETURN" "(" Value ("," Value)* ")" ";"
Binding    ::= Port (":" Type)?
Op         ::= (OpName | Effect) ("." Mode)?
Mode       ::= Ident | ScalarType
Args       ::= Arg ("," Arg)*
Arg        ::= (Ident "=")? (Value | Label)
Value      ::= Port | Ident | Literal
Port       ::= "%" Int
Label      ::= "@" Ident
Literal    ::= Signed ":" ScalarType | "true" | "false" | Array ":" Type
Signed     ::= "-"? Number
Array      ::= "[" Element ("," Element)* "]"
Element    ::= Signed | "true" | "false" | Array
Effects    ::= "[" (Effect ("," Effect)*)? "]"
Effect     ::= "ARENA" | "CLOCK" | "FILE_READ" | "FILE_WRITE" | "NET_READ" | "NET_WRITE"
             | "GPU" | "NONDET" | "PARAM_UPDATE" | "TELEMETRY" | "HOST_CALL" | "SPAWN"
Type       ::= ScalarType | "Bool" | "Eff" | Handle
             | ("Array" | "Tensor") "<" Type "," Shape ">" | "Slice" "<" Type ">"
             | "Tuple" "<" Type ("," Type)* ">" | "Result" "<" Type "," Type ">"
             | "Arena" "<" Int ">" | "{" Ident ":" Type "|" Formula "}"
Shape      ::= Dim | "[" Dim ("," Dim)* "]"
Dim        ::= Int | Ident
Handle     ::= "FileIn" | "FileOut" | "Socket" | "Clock" | "GpuQueue"
ScalarType ::= IntType | FloatType
IntType    ::= "u8" | "u16" | "u32" | "u64" | "i8" | "i16" | "i32" | "i64"
FloatType  ::= "f16" | "bf16" | "f32" | "f64"
Formula    ::= Or ("==>" Formula)?
Or         ::= And ("||" And)*
And        ::= Cmp ("&&" Cmp)*
Cmp        ::= Sum (("==" | "!=" | "<" | "<=" | ">" | ">=") Sum)?
Sum        ::= Prod (("+" | "-") Prod)*
Prod       ::= Unary (("*" | "/" | "%") Unary)*
Unary      ::= ("!" | "-") Unary | Postfix
Postfix    ::= Primary ("[" Formula ("," Formula)* "]" | "." Int)*
Primary    ::= Value | Number | "result" | Call | "(" Formula ")"
             | "forall" "(" Range ("," Range)* ")" "{" Formula "}"
Range      ::= Ident "in" Sum ".." Sum
Call       ::= (("spec" | "fold") ":")? Ident "(" (Formula ("," Formula)*)? ")"
\end{lstlisting}

\section{Reproducing the measurements}\label{app:repro}

All commands run from the repository root at commit \texttt{fe499c8}, on the host described in
\cref{sec:case}.

\paragraph{Tests.} \kw{cargo test --workspace --locked} reports 26 passing tests. The G1
evidence of \cref{tab:g1} is regenerated by \kw{make gate-g1}, which writes
\file{reports/gate-g1.md}; fuzzing takes an hour per target.

\paragraph{The capacity collision (B1).} Generate two modules, each with one node
\kw{NODE @f \{ SIGNATURE: (x: T) -> (T); EFFECTS: []; CONTRACT \{ \} BODY \{ RETURN(x); \} \}},
with $T$ equal to \kw{Tuple<Tuple<u8, ..., u8>, u8>} (257 inner members) and to
\kw{Tuple<Tuple<u8>, u8, ..., u8>} (257 outer members after the inner tuple). Then
\kw{target/debug/vg hash} prints, for both,
\begin{center}\footnotesize
node \texttt{e271d220c54d0cebca0ac788b44471e44d2ffb4039ac865c53e2a301fa7e1c2f}
\end{center}
with equal body and contract hashes, whereas the pair with 2 and 1 inner members prints the
distinct node hashes \texttt{\#25e623e3de21} and \texttt{\#d3edd905070c}.
\kw{target/debug/vg asm} rejects the wide program with \texttt{VG-S0202 unknown type tag 0x00}.

\paragraph{The ordering example.} \Cref{lst:ordering}, preceded by \kw{module ordering;}, is
accepted by \kw{vg asm}, which assigns \kw{@g} and \kw{@f} the node hashes
\texttt{\#78976d2a646b} and \texttt{\#e05de7734ca2}. No verifier exists to run on it; the
proposition is proven by hand.

\paragraph{The soft-knee enumeration.} Compile and run
\begin{center}\footnotesize
\texttt{clang -std=c11 -O2 -ffp-contract=off -fno-fast-math -pthread softknee\_check.c -lm -o
softknee\_check \&\& ./softknee\_check 16}
\end{center}
in \file{docs/paper}. The argument is the number of threads. The program prints the draft's
trace at $x=3\cdot10^{38}$, then the rows E1 to E4 of \cref{tab:enum}; all of its inputs are
fixed or derived from fixed seeds, so the counts are reproducible. The output of the run reported
in \cref{tab:enum} is kept in \file{docs/paper/softknee\_check.out}.
